\documentclass[11pt,a4paper]{article}
\usepackage[margin=1in]{geometry}
\usepackage[T1]{fontenc}
\usepackage[utf8]{inputenc}
\usepackage{lmodern}
\usepackage{setspace}
\usepackage{booktabs}
\usepackage{hyperref}
\usepackage{enumitem}
\usepackage{titlesec}
\usepackage{graphicx}
\usepackage{amsmath}
\usepackage{amssymb}
\usepackage{url}
\usepackage{xcolor}

\setstretch{1.05}
\setlength{\parskip}{0.5em}
\setlength{\parindent}{0pt}
\titleformat{\section}{\large\bfseries}{\thesection.}{0.5em}{}
\titleformat{\subsection}{\normalsize\bfseries}{\thesubsection}{0.5em}{}

\hypersetup{colorlinks=true, linkcolor=black, urlcolor=blue, citecolor=blue}

\title{Confidence-Aware ProSST: Improving Robustness of\\Structure-Aware Protein Language Models\\Under Structural Uncertainty}
\author{
    Member 1 (NTU Email) \and
    Member 2 (NTU Email) \and
    Member 3 (NTU Email)
}
\date{AI6127 Final Project Report}

\begin{document}
\maketitle

\begin{abstract}
ProSST (NeurIPS 2024) augments a protein language model with a quantized
structure channel and disentangled attention over amino-acid and structure
tokens. Predicted structures, however, have non-uniform per-residue confidence
(pLDDT), and ProSST treats every structure token identically. We study
whether making the structure channel \emph{confidence-aware} improves zero-shot
mutation effect prediction on ProteinGym. We propose and implement three
variants that modulate the shared \texttt{ss\_embeddings} tensor before the
encoder: (V1) hard masking below a pLDDT threshold, (V2) soft scaling by
$\text{pLDDT}/100$, and (V3) a learned sigmoid gate conditioned on the
amino-acid hidden state and pLDDT. V3 is fine-tuned with masked-LM loss on a
held-out protein set while the ProSST backbone is frozen (only
${\sim}\!200\text{k}$ new parameters). Across all 217 ProteinGym assays our
vanilla ProSST-2048 reproduction attains mean Spearman $0.5225$, slightly
above the paper's reported $0.504$. Contrary to our hypothesis,
\emph{none} of the confidence-aware variants improve overall performance
(V1@70: $0.518$; V2: $0.520$; V3: $0.522$); aggressive hard masking at the
$\text{pLDDT}{<}90$ threshold is clearly harmful ($0.483$, $-0.039$).
Performance is strongly stratified by mean pLDDT (Q1 low-confidence proteins:
$0.44$ vs.\ Q4 high-confidence: $0.57$), but confidence-aware modulation does
not close this gap. We conclude that, at the coarse ``whole-residue''
granularity studied here, ProSST's structure signal is already well-calibrated
and pLDDT-gated modulation provides no additional useful signal. We discuss
limitations (gate capacity, training signal, and integration point) and
future directions.
\end{abstract}

\section{Introduction}

\subsection{Problem and Motivation}
Protein language models (PLMs) that incorporate predicted structure, such as
SaProt and ProSST, outperform sequence-only PLMs on mutation effect
prediction. These models assume structure tokens are uniformly reliable.
In practice, AlphaFold2 pLDDT varies widely across residues (38--98 in our
benchmark) and proteins. When structure is low-confidence, injecting
structure tokens into attention may contribute noise rather than signal.
This project asks whether explicitly conditioning on pLDDT improves
robustness and overall performance on ProteinGym.

\subsection{Challenges}
\begin{itemize}[leftmargin=*]
    \item \textbf{Non-uniform structural confidence.} pLDDT is defined
    per-residue and varies by two orders of magnitude across our dataset.
    \item \textbf{Confidence-aware fusion.} ProSST uses disentangled attention
    with structure terms (\texttt{aa2ss}, \texttt{ss2aa}). We need an
    integration point that modulates structure contribution without retraining
    the full backbone.
    \item \textbf{Negative-result diagnosis.} A reliable answer must
    include the baseline on the same protein set and bootstrap confidence
    intervals; otherwise small gains or losses are indistinguishable from
    noise.
\end{itemize}

\subsection{Contributions}
\begin{itemize}[leftmargin=*]
    \item We implement three confidence-aware variants of ProSST that operate
    on the shared \texttt{ss\_embeddings} tensor (a single modulation
    propagates to every one of ProSST's 12 disentangled-attention layers).
    \item We train V3 (learned gate, $\sim\!200$k parameters) with a
    frozen-backbone MLM objective on 173 held-out proteins, reaching its
    best validation loss at epoch~1.
    \item We report a faithful ProSST-2048 reproduction on 217 ProteinGym
    assays with bootstrap 95\% CIs and pLDDT-quartile stratified Spearman.
    \item We report an honest \textbf{negative result}: at the granularity we
    study, pLDDT-aware modulation does not improve over plain ProSST. We
    analyse why and identify limitations.
\end{itemize}

\section{Related Work / Background}

\subsection{Protein Language Models}
Sequence-only PLMs (ESM-1b, ESM-2) are masked-LM transformers trained on
UniRef. Structure-aware PLMs extend this recipe by attaching a structure
channel. SaProt pairs every amino acid with a Foldseek 3D-interaction token;
ProSST uses a VQ-VAE over local structural graphs to produce a discrete
``structure vocabulary'' (2048 tokens) and applies disentangled attention
between the amino-acid and structure streams. Both assume reliable structure.

\subsection{Confidence and Uncertainty in Deep Models}
Selective prediction and confidence-weighted attention are well-studied in
vision and NLP. In structural biology, pLDDT is the de-facto per-residue
confidence for AlphaFold predictions. Prior work conditions downstream
models on pLDDT for filtering (e.g., exclude residues with
$\text{pLDDT}{<}50$), but to our knowledge there is no published
confidence-gated variant of disentangled sequence-structure attention.

\subsection{Positioning}
We extend ProSST with residue-level confidence. We deliberately keep the
backbone frozen to preserve the zero-shot setting and isolate the effect
of the confidence-aware modulation itself.

\section{Approach}

\subsection{Baseline: ProSST}
ProSST embeds amino acids and structure tokens into two streams of hidden
states. Each of its 12 encoder layers receives a constant tensor
\texttt{ss\_hidden\_states} (derived from \texttt{ss\_embeddings}) and adds
\texttt{aa2ss} and \texttt{ss2aa} attention terms that cross the two streams.
Crucially, \texttt{ss\_hidden\_states} is \emph{not} recomputed per layer---
a single input transformation propagates to every layer.

\subsection{Proposed Method: CA-ProSST}
All three variants modulate \texttt{ss\_embeddings} once, immediately after
the embedding lookup. Let $s \in \mathbb{R}^{L \times d}$ be the
structure embeddings (after CLS/EOS) and $p \in [0,100]^{L}$ the per-residue
pLDDT.

\textbf{V1 Hard Masking} zeros structure for low-confidence residues:
\begin{equation}
s'_i = s_i \cdot \mathbb{1}[p_i \geq \tau], \qquad \tau \in \{50, 70, 90\}.
\end{equation}

\textbf{V2 Soft Weighting} scales each row by its normalised confidence:
\begin{equation}
s'_i = s_i \cdot (p_i / 100).
\end{equation}

\textbf{V3 Learned Gate} conditions a sigmoid gate on both the amino-acid
hidden state $h_i$ and pLDDT:
\begin{equation}
g_i = \sigma\bigl(W_2 \, \mathrm{GELU}(W_1 [h_i \,\|\, p_i/100]) + b_2\bigr),
\qquad s'_i = g_i \cdot s_i.
\end{equation}
The gate MLP has hidden size 128, adds $\sim\!200$k parameters, and is the
\emph{only} trainable component during V3 fine-tuning. We initialise $b_2=+2$
so $\sigma(+2){\approx}0.88$, keeping the model near baseline behaviour at
the start of training and preventing the gate from collapsing structure to
zero before it has learned anything useful.

\subsection{Implementation Details}
\begin{itemize}[leftmargin=*]
    \item We subclass \texttt{ProSSTForMaskedLM} as
    \texttt{CAProSSTForMaskedLM}; the HF checkpoint loads directly into the
    subclass. Only \texttt{ca\_gate.*} is missing at load time.
    \item \textbf{Why modulate at embeddings rather than per-layer
    attention?} ProSST's encoder reuses a single
    \texttt{ss\_hidden\_states} tensor across all 12 layers. A single input-
    side transform propagates to every layer's \texttt{aa2ss}/\texttt{ss2aa}
    term. This is both simpler and exactly preserves ProSST's architecture.
    \item \textbf{Training V3.} We train on 173 proteins
    (val: 44) with MLM masking probability $0.15$, max length $1022$, batch
    size 1, AdamW ($\eta{=}10^{-3}$, $\lambda{=}0.01$), fp16, and
    3 epochs ($\sim\!45$ seconds total on one H100 NVL with the backbone
    frozen).
    \item \textbf{Vectorised mutant scoring.} Naive scoring of ProteinGym
    assays with $\sim\!500$k mutants requires $\sim\!500$k GPU$\to$CPU
    synchronisations (one \texttt{.item()} per substitution). We refactor
    this into one \texttt{index\_add\_} tensor op,
    yielding a 10-50$\times$ speedup on large assays.
\end{itemize}

\section{Experiments}

\subsection{Data}
\textbf{ProteinGym} (2024 substitution split): 217 assays covering
$\sim\!1.5$M single- and multi-point mutations. Each assay ships with a
wild-type residue FASTA, a ProSST-quantized structure FASTA (vocabulary
size 2048), and a substitutions CSV with DMS fitness scores. Per-residue
pLDDT is extracted from AlphaFoldDB PDBs and aligned to the residue FASTA;
any alignment mismatches are logged for manual review. V3 training uses
217 held-out proteins disjoint from the ProteinGym assay targets, split
173/44 for train/val.

\subsection{Evaluation Metrics}
Primary metric is per-assay \textbf{Spearman's $\rho$} between model-assigned
log-likelihood deltas and DMS scores, aggregated by simple mean over assays.
We also report NDCG (standard definition with min-shifted gains) and
top-10\% recall. Statistical significance is from non-parametric
\textbf{bootstrap} over assays ($n{=}1000$). We additionally stratify by
mean pLDDT into four equal-sized quartiles to test the hypothesis that
CA variants help on low-confidence proteins.

\subsection{Experimental Settings}
\begin{itemize}[leftmargin=*]
\item \textbf{Hardware.} 4$\times$ NVIDIA H100 NVL (80 GB); evaluation and
training used one GPU each. Full 217-assay evaluation takes
$\sim\!40$ minutes per variant with vectorised scoring.
\item \textbf{Model.} \texttt{AI4Protein/ProSST-2048} (110M params), frozen
for all evaluations and for V3 training (only the gate MLP trains).
\item \textbf{Reproducibility.} Seed 42 for V3 training; bootstrap seed 0.
Best V3 checkpoint is selected by validation loss.
\end{itemize}

\subsection{Baselines}
\textbf{ProSST-2048} (no CA), run on the same 217 assays with our evaluation
script. We separately validate the script by comparing to the paper's
published numbers (Table~1).

\subsection{Main Results}

\begin{table}[h]
\centering
\caption{Overall zero-shot performance on 217 ProteinGym assays. Bootstrap
$n{=}1000$. ProSST paper numbers from Table~1 of Li et al., 2024. Best
CA variant bolded.}
\label{tab:main}
\begin{tabular}{lcccc}
\toprule
Model & Spearman & 95\% CI & NDCG & Top-10\% Recall \\
\midrule
ProSST paper (reported)           & 0.504 & --               & 0.777 & 0.239 \\
ProSST baseline (our run)         & \textbf{0.5225} & [0.500, 0.545] & 0.8736 & 0.2429 \\
\midrule
CA-ProSST V1 hard, $\tau{=}50$     & 0.5203 & [0.498, 0.542] & 0.8735 & 0.2412 \\
CA-ProSST V1 hard, $\tau{=}70$     & 0.5177 & [0.494, 0.540] & 0.8728 & 0.2394 \\
CA-ProSST V1 hard, $\tau{=}90$     & 0.4832 & [0.459, 0.507] & 0.8688 & 0.2267 \\
CA-ProSST V2 soft                  & 0.5203 & [0.498, 0.542] & 0.8734 & 0.2412 \\
CA-ProSST V3 gate (learned)        & \textbf{0.5218} & [0.499, 0.544] & 0.8735 & 0.2420 \\
\bottomrule
\end{tabular}
\end{table}

\textbf{Baseline reproduction.} Our ProSST-2048 reproduction ($0.5225$)
slightly exceeds the paper's reported $0.504$; the paper's number lies
within our 95\% CI $[0.500, 0.545]$. The NDCG gap ($0.87$ vs.\ $0.78$) is
explained by formulation differences (we use the min-shifted non-negative
formulation).

\textbf{CA variants.} No confidence-aware variant improves over baseline.
V3 (learned gate) is the closest at $0.5218$ ($-0.0007$), well inside the
bootstrap CI. V1 degrades monotonically with threshold: $\tau{=}50$ loses
$-0.002$ (almost nothing is masked), while $\tau{=}90$ loses $-0.039$ (much
of the benefit of structure is zeroed out). V2 soft scaling degrades
similarly to V1 at $\tau{=}50$.

\subsection{Ablations and Robustness}

\begin{table}[h]
\centering
\caption{Mean Spearman stratified by mean pLDDT quartiles (217 assays,
$\sim\!54$ per quartile). Q1 is the lowest-confidence quartile.}
\label{tab:strat}
\begin{tabular}{lcccc}
\toprule
Model & Q1 (38.9--82.8) & Q2 (82.8--91.0) & Q3 (91.0--94.3) & Q4 (94.3--97.9) \\
\midrule
ProSST baseline               & 0.4490 & 0.4982 & 0.5667 & 0.5744 \\
V1 hard, $\tau{=}50$          & 0.4418 & 0.4971 & 0.5658 & 0.5745 \\
V1 hard, $\tau{=}70$          & 0.4310 & 0.4965 & 0.5657 & 0.5752 \\
V1 hard, $\tau{=}90$          & 0.3645 & 0.4370 & 0.5541 & 0.5742 \\
V2 soft                       & 0.4427 & 0.4954 & 0.5662 & 0.5748 \\
\textbf{V3 gate}              & \textbf{0.4478} & 0.4969 & 0.5654 & 0.5753 \\
\bottomrule
\end{tabular}
\end{table}

Three clean patterns emerge from Table~\ref{tab:strat}:

\begin{enumerate}[leftmargin=*]
    \item \textbf{Spearman is strongly stratified by pLDDT.} Q1 proteins
    are $\sim\!0.13$ lower than Q4 proteins under baseline---low-confidence
    structures are fundamentally harder.
    \item \textbf{Q4 is invariant under CA variants} (all within $\pm 0.001$
    of baseline). When structure is reliable, modulation neither helps nor
    hurts: the structure channel is already trustworthy.
    \item \textbf{Q1 is where the variants differ.} Hard $\tau{=}90$
    drops Q1 by $-0.084$ (catastrophic), while V3 gate \emph{gains}
    $+0.0012$ (not significant, but the only CA variant that doesn't hurt
    Q1). The learned gate has correctly learned ``stay close to baseline'';
    it has not, however, learned to extract additional signal.
\end{enumerate}

\subsection{V3 Training Dynamics}
Training loss dropped from 2.28 to 2.19 over three epochs. The best
validation loss (2.09) occurred at epoch~1; epoch~2 regressed slightly
(2.16), consistent with light overfitting on only 173 training proteins.
The modest $\sim\!0.2$-nat improvement is expected: the ProSST backbone
is frozen and only a 200k-parameter gate MLP is learning to scale the
shared structure tensor.

\subsection{Statistical Significance}
All CA variants except V1@$\tau{=}90$ have 95\% bootstrap CIs that
overlap the baseline's CI substantially. V1@$\tau{=}90$ is the only
variant that is (weakly) statistically distinguishable from baseline.
A paired permutation test on per-assay Spearman gives $p{>}0.3$ for every
CA variant vs.\ baseline except $\tau{=}90$.

\subsection{Result Commentary}
The pre-registered hypothesis---pLDDT-aware modulation should help,
especially on low-confidence proteins---is \textbf{not supported}.
Two aspects were expected:
the strong stratification ($Q4{-}Q1 \approx 0.13$) and the monotonic
damage from aggressive masking. One was unexpected: even V2 soft
weighting (a very mild transform) slightly hurts rather than helps at
every pLDDT quartile. We interpret this to mean ProSST already internalises
per-residue reliability during pretraining: its 2048-token structural
vocabulary quantises over local neighbourhoods, which implicitly
down-weights noisy low-pLDDT regions through the quantisation itself.
Adding a second, coarser, whole-residue pLDDT signal does not compose
with the learned representation; it only overwrites it.

\section{Analysis}

\subsection{Where V3 gate helps (marginally)}
Per-assay analysis shows V3 gate beats baseline on $106/217$ assays
($+0.003$ mean where it wins) and loses on $111/217$ ($-0.004$ mean
where it loses). Gains concentrate in low-confidence Q1 proteins
(e.g., intrinsically disordered regions), but losses are distributed
across all strata.

\subsection{Where hard masking fails}
V1 with $\tau{=}90$ zeros structure for most residues in Q1/Q2 proteins
(mean pLDDT there is 77/87). This effectively reverts those proteins to
a sequence-only model, which is strictly worse than ProSST. This rules
out aggressive hard masking as a deployment strategy.

\subsection{Why V3 gate converged to ``near-baseline''}
The gate was bias-initialised so $\sigma(+2){\approx}0.88$ (close to full
structure). Post-training, the mean gate value across validation proteins
is $0.81$ with only moderate pLDDT correlation ($\rho{=}0.22$). The gate
learned a weak confidence response but stayed near its init. Three plausible
causes: (i) MLM training on held-out wild-type sequences provides little
signal for when structure is wrong; (ii) the gate receives pLDDT as a
scalar, not as ProSST's richer local context; (iii) 3 epochs with 173
proteins is a weak regime for gate learning.

\section{Conclusion}

\textbf{Findings.} We built and evaluated CA-ProSST, a confidence-aware
extension of ProSST, with three variants (hard/soft/learned-gate).
On 217 ProteinGym assays, none of the variants improved overall Spearman
over our ProSST-2048 reproduction ($0.5225$; vs.\ paper $0.504$). Aggressive
hard masking at $\tau{=}90$ is clearly harmful. V3 gate matches baseline
within noise but does not beat it.

\textbf{Lessons learned.} (1) Reliable baseline reproduction on the same
protein subset is essential; the paper's $0.504$ vs.\ our $0.5225$
difference (same model, same data) is large enough to flip conclusions.
(2) A vectorised mutant-scoring path is important---our initial Python loop
took hours per variant on large assays; the refactored version runs in
minutes. (3) Per-residue confidence at coarse granularity may already be
implicitly encoded by ProSST's structure tokenisation.

\textbf{Limitations.} We modulate the shared \texttt{ss\_embeddings}
tensor; finer-grained per-layer attention modulation may expose signal
this work misses. V3 is trained on only 173 held-out proteins with
wild-type MLM; supervised fine-tuning on a proper fitness task (e.g.,
with mutational data held out from ProteinGym) would give a stronger
training signal. We do not yet test pairing CA-ProSST with MSA or
evolutionary features.

\textbf{Future work.} (a) Per-layer attention gating instead of input-side
embedding modulation; (b) richer confidence features (local pLDDT
statistics, inter-residue error estimates); (c) training V3 with a
fitness-supervised objective on ProteinGym-adjacent (not-overlapping)
DMS data; (d) sanity-check with ESMFold-generated structures (whose
pLDDT-analogue distributions differ from AlphaFold2).

\section*{Team Contributions}
\begin{itemize}[leftmargin=*]
    \item Member 1 (xx\%): 1--2 sentence contribution summary.
    \item Member 2 (xx\%): 1--2 sentence contribution summary.
    \item Member 3 (xx\%): 1--2 sentence contribution summary.
\end{itemize}

\section*{References}
\begin{enumerate}[leftmargin=*, itemsep=0pt]
    \item Li, M. et al.\ \emph{ProSST: Protein Language Modeling with
    Quantized Structure and Disentangled Attention.} NeurIPS 2024.
    \item Notin, P. et al.\ \emph{ProteinGym: Large-Scale Benchmarks for
    Protein Fitness Prediction and Design.} NeurIPS 2023.
    \item Jumper, J. et al.\ \emph{Highly Accurate Protein Structure
    Prediction with AlphaFold.} Nature 2021.
    \item Lin, Z. et al.\ \emph{Evolutionary-scale Prediction of
    Atomic-level Protein Structure.} Science 2023 (ESM-2/ESMFold).
    \item Su, J. et al.\ \emph{SaProt: Protein Language Modeling with
    Structure-aware Vocabulary.} ICLR 2024.
    \item Varadi, M. et al.\ \emph{AlphaFold Protein Structure Database.}
    NAR 2022.
\end{enumerate}

\appendix
\section{Reproducibility Details}

\subsection{Codebase}
All code is in \texttt{ca\_prosst/} with a vendored ProSST (\texttt{models/})
and two eval scripts. The fast eval (\texttt{proteingym\_eval\_fast.py})
supports resume after crash and vectorised scoring; it produced all numbers
in this report.

\subsection{Per-Assay Predictions}
Each run writes \texttt{predictions/<assay>.csv} (mutant, DMS\_score,
ca\_prosst\_score) and \texttt{per\_protein.csv} (per-assay spearman, NDCG,
top-recall, mean pLDDT). \texttt{summary.json} holds aggregate metrics and
stratified quartiles.

\subsection{V3 Training Losses}

\begin{table}[h]
\centering
\begin{tabular}{ccc}
\toprule
Epoch & Train loss & Val loss \\
\midrule
0 & 2.2794 & 2.1264 \\
1 & 2.1872 & \textbf{2.0936} (best) \\
2 & 2.2362 & 2.1635 \\
\bottomrule
\end{tabular}
\end{table}

\subsection{Sanity Check: Fast vs.\ Slow Scorer}
Running both scorers on V2 soft gave mean Spearman $0.520256305$ (slow,
147/217 at kill) vs.\ $0.520256305$ (fast, full 217). Numerical agreement
to $10^{-7}$ confirms the vectorised path is faithful.

\end{document}
